\subsection{Conclusions\label{sec:conc}}
This study presents the first simultaneous time-resolved measurements of key gas-phase species and soot volume fraction during shock-tube pyrolysis of 5\%, 10\%, and 30\% methane over a temperature range of 1800–2500~K and pressures near 4~atm. Using multi-wavelength laser diagnostics, we quantified the evolution of CH$_4$, C$_2$H$_4$, and C$_2$H$_2$. The results show that methane conversion and intermediate yields strongly depend on temperature and fuel concentration, with significant model–measurement deviations for CH$_4$ mole fraction emerging at 30\% loading due to increased absorption interference. Comparisons with detailed kinetic models revealed that the FFCM2 and CRECK models capture gas-phase species more accurately despite significant differences in pathways. Carbon flux analysis highlighted the dominant role of C$_3$H$_3$ in channeling carbon toward C$_2$H$_2$ and PAHs, and the role of P-C$_3$H$_4$ as a key precursor of C$_3$H$_3$. The model captured the time and temperature trends of $f_v$ at 5\% and 10\% loadings, but a substantial overprediction with similarities in time trends is observed at 30\% loading, indicating potential missing chemistry or physical effects that could inhibit surface growth at higher loadings.